\section*{What this book is}

\emph{Idea Theory} is the proposal that there is a single mathematical
structure underlying the very many things human beings have wanted to call
``ideas'' --- the noemata of phenomenology, the units of conceptual
history, the strategies of evolutionary game theory, the tokens of cultural
transmission, the latent vectors of a transformer, the categories of Kant,
the Aufhebung of Hegel, the memes of Dawkins, the nonrival goods of Romer,
and the emergent qualities of the British emergentists. The structure is
modest --- a graded monoid equipped with a bilinear pairing, a triadic
decomposition of composition, and a 2-cocycle measuring the failure of
valuations to be additive --- yet it suffices to formalise an astonishing
fraction of what those traditions assert.

The book proceeds in two interleaved registers. The first is the
\emph{formalism}: the nine foundational theorems, machine-checked in
Lean~4 under the namespace \texttt{IdeaTheory}, are stated and proved here
as ordinary mathematics, with no Lean code in the body. The second is the
\emph{zoo}: each chapter populates the formalism with concrete entries
from the literature on ideas. Plato's Forms reappear as fixed points of
composition, Hegel's Aufhebung as the triadic decomposition itself,
Hume's association as a bilinear pairing, Anderson's ``more is different''
as a non-trivial cocycle class, Romer's nonrival ideas as a combinatorial
closure, Mikolov's word embeddings as a learned resonance form, and so
on. The intent is not to claim that the formal restatement \emph{exhausts}
any of these positions --- it does not --- but to expose the shared
algebraic skeleton that makes them mutually translatable.

\section*{How to read it}

Chapters \ref{ch:monoid}--\ref{ch:grading} develop the formalism. They are
ordered so that each new piece of structure is introduced, axiomatised,
and worked out in mathematical detail before its zoo is opened. A reader
willing to take the foundational theorems on faith may skim the
math-prose proofs and concentrate on the zoo entries, which are
self-contained.

Chapters~\ref{ch:philosophy}--\ref{ch:econ} are dedicated zoos: the
philosophy of ideas, the cognitive science of concepts, and the
game-theoretic and economic theories of innovation are each reformalised
as a long sequence of zoo entries. These chapters cite the formalism
heavily but do not introduce new mathematical structure.

The appendix lists every Lean module on which the book depends, with a
one-paragraph description. The bibliography lists real, verifiable
publications. Every claim attributed to a thinker is cited.

\section*{What is new}

The novelty is not in any individual theorem --- monoids, bilinear pairings,
group cohomology, and cocycle conditions are old --- but in the claim that
\emph{this particular bundle of structure}, with these particular
compatibility axioms, suffices to translate so much of the literature on
ideas into a single algebraic object. The book exists to test that claim.
We invite the reader to look for a position in their own area that
\emph{cannot} be reformalised within the framework. Where such cases
arise, they are the most interesting.
